\documentclass{article}
\usepackage{amsmath, amssymb, amsthm}
\usepackage{hyperref}
\usepackage{geometry}
\geometry{margin=1in}

\title{Erd\H{o}s Problem \#730}
\date{Last updated: 2025-08-31}
\author{Source: erdosproblems.com}

\begin{document}
\maketitle

\noindent\textbf{Status:} OPEN \quad \textbf{No prize}

\noindent\textbf{Tags:} number theory, binomial coefficients, base representations

\medskip
\noindent\textbf{Problem Statement:}

Are there infinitely many pairs of integers $n\neq m$ such that $\binom{2n}{n}$ and $\binom{2m}{m}$ have the same set of prime divisors?

\bigskip
\noindent\textbf{Remarks:}

A problem of Erd\H{o}s, Graham, Ruzsa, and Straus \cite{EGRS75}, who believed there is 'no doubt' that the answer is yes. 

For example $(87,88)$ and $(607,608)$. Those $n$ such that there exists some suitable $m>n$ are listed as A129515 in the OEIS. 

A triple of such $n$ for which $\binom{2n}{n}$ all share the same set of prime divisors is $(10003,10004,10005)$. It is not known whether there are such pairs of the shape $(n,n+k)$ for every $k\geq 1$.

\bigskip
\noindent\textbf{References:}

[EGRS75] Erd\H{o}s, P. and Graham, R. L. and Ruzsa, I. Z. and Straus, E. G., On the prime factors of $(\sp{2n}\sb{n})$. Math. Comp. (1975), 83-92.

\bigskip
\noindent\small{Source: \url{https://www.erdosproblems.com/730}}
\end{document}
